\scanpage{15}
$\deleted{4}\unsure{3|4})\quad L_i\notin H_{i+1}\qquad 0\leq i\leq n-1$, en tte fibre

i.e. $L_i\otimes L'_{i+1}\xrightarrow{\sim}\underline O_S$ \inserted{hom can. induit par $\deleted{\unsure{V}}V\otimes V^\vee\to\underline O_S$}

$\deleted{\unsure{C'_0}}$

\deleted{$C'_0={}$ Cat des triples précédents \uline{rigidifiés} par une section $a_0$}

\uline{NB} De 4), on tire
\[
\left\{\begin{aligned}
&L_i,L'_{i+1}\text{ en dualité}\qquad (0\leq i\leq n-1)\quad\text{donc}\\
&L_i\xrightarrow{\sim}L_{i+1}'{}^\vee\qquad (0\leq i\leq n-1)
\end{aligned}\right.
\]
De 1) on tire
\[
\left.\begin{aligned}
L'_i\otimes L_i^{\deleted{\unsure{\vee}}}&\xleftarrow{\sim}\underline O_S\\
\tau_i^V-\mathrm{id}_V&\longleftarrow\!\shortmid\,1
\end{aligned}\ \right|\quad 1\leq i\leq n\quad \deleted{\text{d'où}}\quad\text{donc}
\]
\[
L'_i\underset{\mathrm{can}}{\simeq} L_i^{\prime\vee}\qquad 1\leq i\leq n
\]
On en conclut de proche en proche des isos
\[
L_0\simeq L_1^{\prime\vee}\simeq L_1\simeq L_2^{\prime\vee}\simeq L_2\ldots\simeq L_n^{\prime\vee}\simeq L_n
\]
d'où $\backslash$
\[
V\simeq L\otimes\underline O_S^{n+1}\qquad\text{iso can.}
\]

$\boxed{C'_0}={}$ cat des triples $(V,\varphi_G^V,L_0)$ \uline{rigidifiés} par une section\inserted{-base} $a_0$ de $L_0$, i.e. des triples $(V,\varphi_G^V,a_0)$, \deleted{$a_0$ se} $V,\varphi_G^V$ comme dessus, $a_0$ section $\ne0$ en tte fibre, satisfaisant les conditions 1) à 4) où on fait $L_0=\underline O_S a_0$ [Le fait que $a_0$ est $\ne0$ sur chaque fibre est déjà \unsure{contenu} dans la première condition 4), si on l'écrit $a_0\notin H_1$ \deleted{\unsure{en}} tte fibre \ldots]

% archive sheet number: 14

